\begin{table}[t]
\centering\small
\setlength{\tabcolsep}{3pt}
\begin{tabular}{@{}llrrl@{}}
\toprule
& & $\Delta$\,train & $\Delta$\,held & \textbf{retention} \\
& & \multicolumn{1}{c}{\footnotesize gpt-4o-mini} & \multicolumn{1}{c}{\footnotesize Haiku} & \footnotesize [95\% CI] \\
\midrule
\multirow{4}{*}{\textbf{Q1}}
 & PTO\,@10  & 1.22 & 0.97 & \textbf{0.80} [0.68, 0.93] \\
 & GRPO\,@10 & 0.68 & 0.19 & \textbf{0.28} [0.06, 0.43] \\
 & PTO\,@8   & 1.18 & 1.05 & 0.89 [0.76, 1.04] \\
 & GRPO\,@8  & 1.01 & 0.65 & 0.64 [0.52, 0.78] \\
\midrule
\multirow{4}{*}{\textbf{Q2}}
 & PTO\,@10  & 1.30 & 1.11 & 0.85 [0.74, 0.98] \\
 & GRPO\,@10 & 0.82 & 0.66 & 0.81 [0.68, 0.96] \\
 & PTO\,@8   & 1.26 & 1.08 & 0.86 [0.74, 0.99] \\
 & GRPO\,@8  & 1.15 & 0.92 & 0.80 [0.68, 0.95] \\
\bottomrule
\end{tabular}
\caption{Gain retention $=\Delta_\text{held-out}/\Delta_\text{trained-against}$, where
$\Delta$ is the persona-paired improvement over the base policy ($n=96$). \textbf{Q1 separates
the arms with non-overlapping intervals at the matched iteration-10 endpoint}; at GRPO's
iteration-8 peak PTO still leads but the intervals overlap marginally, so the peak-vs-peak
claim is an ordering rather than a disjoint contrast (\S\ref{sec:heldout}). \textbf{Q2 does not
separate at all} --- all four intervals sit in a $0.80$--$0.86$ band, the control showing that
retention is not merely a stingier grader.}
\label{tab:retention}
\end{table}
